
%\needspace{4\baselineskip}
\color{uniblue}{\subsection{Контроль отсутствия напряжения на присоединении (КОНп)}}
\color{black}

\begin{enumerate}[label=\arabic{section}.\arabic{subsection}.\arabic*, labelsep=4pt, leftmargin=0pt, itemindent=57pt]

\item
Функция контролирует отсутствие напряжения на отходящем присоединении. Функция используется для АПВ, а также для оперативного включения.
\item
Контроль отсутствия напряжения на присоединении срабатывает без выдержки времени при одновременном снижении всех линейных напряжений ниже уставки. При срабатывании ИО напряжения ниже уставки, на выходе будет отображаться сигнал «КОНп: ИО U<».
\item
Пусковые условия функции формируются при одновременном наличии следующих условий:
\begin{itemize}
    \item уставка SGF2 выставлена в положение <<Измеренный>>, то при снижении максимального из линейных напряжений присоединения ниже уставки (наличие выходного сигнала «КОНп: ИО U<»), наличие выходного сигнала <<КОНп: Ввод>> и отсутствие входного сигнала срабатывания от неисправности цепей напряжения присоединения <<БННп: Срабатывание>>;
    \item уставка SGF2 выставлена в положение <<Внешний>>, наличие выходного сигнала <<КОНп: Ввод>> и наличие активного входного дискретного сигнала отсутствия напряжения на присоединении <<Отсутствие Uп>>.
\end{itemize}

\item
Блокировка пусковых условий формируется при наличии неисправности цепей напряжения в ТН (БННп) и повышение напряжения ОП выше уставки если уставки SGF3 и SGF2 выставлены в положение <<0>>. Информация о срабатывании ИО ОП будет отображаться выходным сигналом «КОНп: ИО U<».

\item
Предусмотрена блокировка функции при срабатывании ИО обратной последовательности напряжения выше уставки и наличии сигнала срабатывания БННп.

\item
Ввод функции в работу на этапе параметрирования устройства осуществляется программным переключателем SGF1 <<Ввод\_функции>>. Информация о введенной функции отображается выходным сигналом <<КОНп: Ввод>>.
\item
Функция <<КОНп>> может быть выведена из работы оперативно путем активации сигнала <<ОВ КОНп>>, что характеризуется активным состоянием сигнала <<КОНп: Оперативный вывод>> на выходе функции. 
\item
Алгоритм работы функции <<КОНп>> выполнен в соответствии с рисунком \ref{konp:img1}. В таблице \ref{konp:tbl1} приведены параметры, необходимые для настройки функции <<КОНп>>.

\vspace{3mm}
\begin{figure}[H]
\centering
\includegraphics[width=1\textwidth,height=1\textheight,keepaspectratio]{\fbpath/125. КОНп/_latex/img/konp.pdf}
\captionof{figure}{Функциональная схема <<КОНп>>}\label{konp:img1}
\end{figure}

\small
\begin{longtable}{|>{\centering\arraybackslash}m{5.3cm}|>{\centering\arraybackslash}m{3.3cm}|>{\centering\arraybackslash}m{4.2cm}|>{\centering\arraybackslash}m{1.8cm}|>{\centering\arraybackslash}m{1cm}|}
\caption{Параметры для настройки функции <<КОНп>>\hfill\vspace{-0.5\baselineskip}}\label{konp:tbl1}\\ 
\hline
\rowcolor{gray!30}
Параметр (Параметр на ИЧМ) & Условное обозначение на схеме & Значение/ Диапазон & Единица измерения & Шаг \\ 
\hline
\endfirsthead
\caption*{\hspace{3pt}\emph{Продолжение таблицы \ref{konp:tbl1}\hfill\vspace{-0.5\baselineskip}}} \\ % сделано по ГОСТ 2.105 п.6.8.7
\hline
\rowcolor{gray!30}
Параметр (Параметр на ИЧМ) & Условное обозначение на схеме & Значение/ Диапазон & Единица измерения & Шаг \\ 
%\hline
\endhead
%\hline
\endfoot
%\hline
\endlastfoot
%==+t1*PTUV1|VOLABS2
\centering Ввод функции в работу (Ввод\_функции) & \centering SGF1 & \centering 0 = Не предусмотрено\\1 = Предусмотрено & \centering -- & \centering \arraybackslash -- \\
\hline
\centering Напряжение срабатывания (Uср) & \centering U< & \centering 2,0 ... 100,0 & \centering В & \centering \arraybackslash 0,1 \\
\hline
\centering Напряжение срабатывания ОП (U2ср) & \centering U2> & \centering 2,0 ... 50,0 & \centering В & \centering \arraybackslash 0,1 \\
\hline
\centering Контроль ТН (Контроль\_ТН) & \centering SGF3 & \centering 0 = 3 фазы\\1 = 2 фазы & \centering -- & \centering \arraybackslash -- \\
\hline
\centering Режим пуска (Реж\_пуска) & \centering SGF2  & \centering 0 = Измеренный\\1 = Внешний & \centering -- & \centering \arraybackslash -- \\
\hline
%===t1
\end{longtable}
\normalsize

\end{enumerate}
